\section{Conclusion}
\label{sec:conclusion}

ST deconvolution is the masked-data identifiability problem from
reliability statistics, with cell types playing the role of
components, spots playing the role of candidate sets, and
single-cell-resolution platforms playing the role of singleton
candidate sets. We proved (i) a rank condition on the spot-by-cell-
type composition matrix that determines when deconvolution is
identifiable, (ii) a cell-total consistency theorem explaining why
naive deconvolution methods produce plausible aggregate predictions
even when per-cell-type estimates disagree, and (iii) a first-order
bias bound under marker-gene undercoverage that recovers Tangram's
empirical ``a few hundred markers'' rule as a corollary.

The framework subsumes Cell2location, RCTD, Tangram, and CARD as
special cases and provides a precise vocabulary for when each
method's assumptions hold. Practitioners gain a framework-level
diagnostic for method choice; methodologists gain a unified
theoretical scaffold for further work.

The principal open extensions are joint composition--expression
inference, spatial-autocorrelation priors, and incorporation of
single-cell-resolution measurement noise via C2-relaxed coarsening
\cite{towell2026mdrelax}. Each is a natural follow-up.

Spatial transcriptomics is one of several distinct domains where
the masked-data framework yields identifiability conditions and
bias bounds. Companion papers carry the same apparatus to scRNA-seq
zero inflation \cite{towell2026scrnacoarsening}, differential
privacy \cite{towell2026dpcoarsening}, programmatic weak
supervision \cite{towell2026weaksupcoarsening}, and EHR phenotyping
\cite{towell2026phenotypecoarsening}. The portability across these
applications is what motivates the framework-level investment.
